% !PRINTED 563
\weberpage{563}
%<<CONT
zahl der Punkte mit ganzzahligen Coordinaten, die in einem im
Verhältniss $t^{\frac1n} : 1$ vergrösserten Gebiet $G_x^{(t)}$
liegen, und diese Zahl ist ebenso gross wie die Anzahl der
Gitterpunkte, die in dem Gebiete $G_\xi^{(t)}$ liegen. Wir haben
also folgenden Satz:

\gesperrt{Ist $G_\xi$ irgend ein endliches Gebiet im Raume $R_\xi$
und $G_\xi^{(t)}$ das aus $G_\xi$ durch Vergrösserung im
Linearverhältniss $t^{\frac1n} : 1$ abgeleitete Gebiet; ist $Z_t$
die Anzahl der im Inneren von $G_\xi^{(t)}$ liegenden Gitterpunkte,
so ist das durch $G_\xi$ bestimmte Volumen}

\weq{15}{V_\xi = \sqrt{D}\ \mathrm{Lim}_{t\,=\,\infty} \frac{Z_t}{t}.}

Wir legen jetzt um jeden Gitterpunkt als Mittelpunkt eine
$n$-dimensionale Kugel vom Radius $\frac12\sqrt{M}$. Da $\sqrt{M}$
die kleinste Entfernung von irgend zwei Gitterpunkten ist, so können
je zwei dieser Kugeln kein Gebiet gemein haben; sie können sich
höchstens in einem Punkte berühren, und keine dieser Kugeln enthält
ausser ihrem Mittelpunkte einen weiteren Gitterpunkt. Eine solche
Kugel hat nach (10) das Volumen

\[
K_n \left(\frac{\sqrt{M}}{2}\right)^{n},
\]

und das Gesammtvolumen aller der Kugeln, deren Mittelpunkt im
Inneren von $G_\xi^{(t)}$ liegt, ist

\[
V_1^{(t)} = Z_t\, K_n \left(\frac{\sqrt{M}}{2}\right)^{n}.
\]

Dies Volumen $V_1^{(t)}$ ist aber offenbar kleiner als $V^{(t)}$,
wenigstens wenn wir von den Kugeln, die zum Theil aus dem Gebiete
$G^{(t)}$ herausragen, absehen, deren Gesammtvolumen, durch $t$
dividirt, mit unendlich wachsendem $t$ gegen Null convergirt, weil
sein Gebiet nur einen Theil von $G^{(t)}$ erfüllt, und folglich ist
auch $V$ grösser als der Grenzwerth von $V_1^{(t)} : t$, also

\[
V > \mathrm{Lim}\, \frac{Z_t\, K_n}{t} \left(\frac{\sqrt{M}}{2}\right)^{n},
\]

und daraus nach (15)

\weq{16}{\sqrt{D} > K_n \left(\frac{\sqrt{M}}{2}\right)^{n}.}

Da es nun leicht ist, das Volumen $K_n$ zu bestimmen, so ist
hierdurch die Aufgabe gelöst, und eine obere Grenze für $M$
gefunden.
